% Source-derived chapter 05: Multiplicities in the nilpotent cone
% Original source: very-stable-higgs-bundles-equivariant-multiplicity-mirror-symmetry
\chapter{Multiplicities in the nilpotent cone}
\label{very-stable-higgs-bundles-equivariant-multiplicity-mirror-symmetry-chapter-05}
\source{very-stable-higgs-bundles-equivariant-multiplicity-mirror-symmetry}

\begin{remark}[Source section]
\label{very-stable-higgs-bundles-equivariant-multiplicity-mirror-symmetry-chapter-05-source}
\source{very-stable-higgs-bundles-equivariant-multiplicity-mirror-symmetry}
\source{very-stable-higgs-bundles-equivariant-multiplicity-mirror-symmetry}
This chapter reproduces the corresponding section of the retrieved
LaTeX source, including its definitions, statements, examples, and
proofs. It has not yet been translated into Lean.
\notready
\end{remark}

% The source section title was: Multiplicities in the nilpotent cone
\label{multiplicities}
\source{very-stable-higgs-bundles-equivariant-multiplicity-mirror-symmetry}

\subsection{Multiplicities of very stable components}
Let $\calE\in \M^{s \T}\subset h^{-1}(0)$ be a very stable Higgs bundle. Denote by $F_\calE\in \pi_{0}(\M^{ \T})$ the component of the $\T$-fixed point set containing $\calE$.   Let $N:=h^{-1}(0)\subset \M$ denote the subscheme of nilpotent Higgs bundles, the so called {\em (global) nilpotent cone}. Then the corresponding subvariety $N_{red}\cong \calC\subset \M$ is isomorphic to the core of $\M$. Denote by $\calC_{F_{\calE}}:= \overline{W^-_{F_\epsilon}}$ the closure of the downward flow from $F_\calE$, which is an irreducible component of $\calC$. Denote by $N_{F_{\calE}}\subset N$ the corresponding irreducible component of the subscheme $N$.  

\begin{definition}
\source{very-stable-higgs-bundles-equivariant-multiplicity-mirror-symmetry}
\notready \label{defmult} The {\em multiplicity} of the component $N_{F_\calE}$   in the subscheme $N$ is defined to be the length of the local ring $\calO_{N,N_{{F_\calE}}}$ at a generic point  and is  denoted by $m_{F_\calE}:=\ell(\calO_{N,N_{{F_\calE}}})$.   
\source{very-stable-higgs-bundles-equivariant-multiplicity-mirror-symmetry}
	%	Sometimes, to simplify notation we will just write $m_\calE$ for $m_{F_\calE}$.
\end{definition}


 Let $k\in \Z$ be an integer. For a $\T$-module $U=\bigoplus_{\lambda\leq k} U_\lambda$  with finite dimensional weight spaces $\dim(U_\lambda)<\infty$ corresponding to $\lambda\in \Hom(\T,\T)\cong \Z$   we denote its character by \beq\label{trunchar}\chi_\T(U):=\sum_{\lambda\leq k} \dim(U_\lambda)t^{-\lambda} \in \Z((t)).\eeq For a finite dimensional positive $\T$-module $V=\oplus_{\lambda>0} V_\lambda$  the  character of the  symmetric algebra $\Sym(V^*)=\bigoplus_{\lambda\leq 0} \Sym(V^*)_\lambda$ of its dual satisfies
\source{very-stable-higgs-bundles-equivariant-multiplicity-mirror-symmetry}
\beq\label{sym} \chi_\T({\Sym(V^*)})&=& \sum_{\lambda\leq 0} \dim(\Sym(V^*)_\lambda)t^{-\lambda}= % \sum_{\lambda\geq 0} \dim(\Sym(V)_\lambda)t^\lambda\\ &=&
\source{very-stable-higgs-bundles-equivariant-multiplicity-mirror-symmetry}
\prod_{\lambda>0}\frac{1}{( 1- t^{\lambda})^{\dim(V_\lambda)}} \in \Z[[t]].\eeq In particular,  the equivariant Euler characteristic (see \S\ref{equivarianteuler} for more details) of the structure sheaf $\calO_V$ on $V$ can be computed as
\beq \label{eqeu} \chi_\T(V;\calO_V)=\sum_{i,k} (-1)^i t^{-k} \dim H^i(\calO_V)_{k} =\sum_k t^{-k} \dim H^0(\calO_V)_{k} = \chi_\T({\Sym(V^*)}),\eeq where we recall that $\T$ acts on $s
\source{very-stable-higgs-bundles-equivariant-multiplicity-mirror-symmetry}
\in H^0(\calO_V)$ with non-positive weights by the formula $(\lambda\cdot s)(v)=s(\lambda^{-1}\cdot v)$. 

Finally, we let $T_\calE^+\M<T_\calE\M$ denote the part of the $\T$-module $T_\calE\M$ with positive weights. Then $T_\calE W^+_\calE\cong T_\calE^+\M$ which  we will abbreviate as $T_\calE^+$.
\begin{theorem}
\source{very-stable-higgs-bundles-equivariant-multiplicity-mirror-symmetry}
\notready \label{multiplicity} When $\calE\in \M^{s \T}\subset N$ is a very stable Higgs bundle then the multiplicity of the component $N_{F_\calE}$   in the nilpotent cone $N$ satisfies $$m_{F_\calE}=\rank(h_*(\calO_{W^+_\calE}))=\left. { \chi_\T(\Sym(T^{+*}_\calE))\over \chi_\T(\Sym(\calA^*))}\right|_{t=1}.$$  Moreover, for a generic $a\in \calA$  the intersection $W^+_\calE\cap h^{-1}(a)$ is transversal and has cardinality $m_{F_\calE}$. 
\source{very-stable-higgs-bundles-equivariant-multiplicity-mirror-symmetry}
\end{theorem}
\begin{proof} We shall use basic properties of the 
	K-theory of coherent sheaves with supports as developed in e.g. \cite{baum-etal,gillet-soule, piepmeyer-walker}. For $Z\subset X$ a closed algebraic  subset of a complex algebraic variety we will denote by $K^\circ_Z(X)$ (resp. by $K^Z_\circ(X)$)  the Grothendieck group of bounded complexes of locally free sheaves (resp. coherent sheaves) with homology supported in $Z$.
	
	We shall compute the intersection product \beq \label{rankmult} [\calO_{N}]\cap [\calO_{W^+_\calE}]\in  K_\circ(\{\calE\})\cong \Z\eeq 
\source{very-stable-higgs-bundles-equivariant-multiplicity-mirror-symmetry}
	in two different ways. Here $[\calO_{N}]=h^*([\calO_0])\in K^\circ_\calC(\M)$, where  $[\calO_0]\in K^\circ_{\{0\}}(\calA)$
	and $[\calO_{W^+_\calE}]\in K^{W^+_\calE}_\circ(\M)$, where $W^+_\calE\subset \M$ is closed by Lemma~\ref{transversal}.  Thus the intersection product $$[\calO_{N}]\cap[\calO_{W^+_\calE}]=[\calO_N\otimes^L\calO_{W^+_\calE}] \in K_\circ^{\calC\cap W^+_{\calE}}(\M)\cong K^{\{\calE\}}_\circ (\M)\cong K_\circ(\{\calE\})\cong \Z$$ is defined. The last isomorphism is given by the Euler characteristic $\chi:K_\circ(\{ \calE \})\to \Z$.  % Note that $h^*: K^\circ_c(\calA) \to K^\circ_c(\M)$ is well-defined as the Hitchin map is proper. 
	%As the support of the complex $\calO_{W^+_\calE}\otimes^L \calO_{N}$ is $\{\calE\}\subset \M$ we have $$[\calO_{W^+_\calE}]\otimes [\calO_{N}]=[\calO_{W^+_\calE}\otimes^L \calO_{N}]=[\calO_{W^+_\calE}\otimes^L \calO_{N_{F_\calE}}]\in K_\circ^c(\M).$$
	
	
	First we note that $W^+_\calE$ intersects $\calC=N_{red}$ transversally at the single point $\calE$, since both $W^+_\calE$ and $\calC$ are smooth at $\calE$ and the tangent spaces  $T^+_\calE$ and $T^{\leq  0}_\calE$ are complementary. This also follows from Bialynicki-Birula's local description of the upward flow (see Remark~\ref{afffibr}). Then \cite[\href{https://stacks.math.columbia.edu/tag/0B1I}{Lemma 42.14.3}]{stackproject} implies that the intersection multiplicity \beq\label{tranint}i(\M,W^+_\calE\cdot \calC,\{\calE\})=i(\M,W^+_\calE\cdot \calC_{F_\calE},\{\calE\})=1.\eeq
\source{very-stable-higgs-bundles-equivariant-multiplicity-mirror-symmetry}
	
	We can also find an affine open neighbourhood $\spec(R)$  of $\calE\in \M$ so that $\calO_{W^+_\calE}|_{\spec(R)}$ is represented by the $R$-algebra $A$ and $\calO_{N_{F_\calE}}$ by the $R$-algebra $B$. As ${N_{F_\calE}}$ is irreducible $B$ has a unique minimal prime ideal $I_{min}\lhd B$ which corresponds to the generic point of ${N_{F_\calE}}$ and coincides with the nilradical of $B$. In particular, the reduced ring $B_{red}=B/I_{min}$ and \beq\label{reduced}\spec (B_{red}) \cong \calC_{F_{\calE}}\cap{\spec(R)}.\eeq  By \cite[Theorem 2 of IV.1.4]{bourbaki} %\footnote{We thank Quoc Ho for suggesting to use this result.} 
\source{very-stable-higgs-bundles-equivariant-multiplicity-mirror-symmetry}
	we know that $B$ as a $B$-module has a composition series $$0=B_k\subset B_{k-1} \subset \dots\subset B_0=B$$ with factors $B_j/{B_{j+1}}\cong B/P_j$, where $P_j\lhd B$ are prime ideals. Moreover we know from  {\em loc. cit.} that the factor $B/I_{min}$ appears $\ell(B_{I_{min}})$ times in this composition series, which further agrees with $$\ell(B_{I_{min}})=\ell(\calO_{N,N_{{F_\calE}}})=m_{F_\calE}.$$ 
	
	Now $\calO_{W^+_\calE}\otimes^L \calO_{N_{F_\calE}}$ has support $\{\calE\}$. The Euler characteristic of its stalk at $\{\calE\}$ is given by Serre's Tor formula
	$$\chi(A_{I_\calE},B_{I_\calE})=\sum (-1)^i \ell(\Tor_i(A_{I_\calE},B_{I_\calE})),$$
	where $I_\calE\lhd R$ is the maximal ideal corresponding to the point $\calE\in \spec(R)\subset \M$. Computing this further
	$$\sum_i (-1)^i \ell(\Tor_i(A_{I_\calE},B_{I_\calE}))=\sum_i \sum_j (-1)^i \ell(\Tor_i(A_{I_\calE},(B/P_j)_{I_\calE}))=\sum_j \chi(A_{I_\calE},(B/P_j)_{I_\calE}).$$
	We know that $P_j\supset I_{min}$. When $P_j\neq I_{min}$ then $\dim(B/P_j)<\dim(B/I_{min})$ and so Serre's \cite[Theorem C.1.1.a]{serre} shows that $\chi(A_{I_\calE},(B/P_j)_{I_\calE})=0$ in this case. Thus  the Euler characteristic of the stalk of $\calO_{W^+_\calE}\otimes^L \calO_{N_{F_\calE}}$ at $\{\calE\}$ equals $$\chi(A_{I_\calE},B_{I_\calE})=m_{F_\calE} \chi(A_{I_\calE},(B/I_{min})_{I_\calE})=m_{F_\calE},$$
	as $$\chi(A_{I_\calE},(B/I_{min})_{I_\calE})=\chi(A_{I_\calE},(B_{red})_{I_\calE})=i(\M,W^+_\calE\cdot \calC_{F_\calE},\{\calE\})=1$$ from \eqref{tranint},\eqref{reduced} and \cite[Theorem C.1.1.b]{serre}. This  implies  that \beq \label{intermult}\chi([\calO_{W^+_\calE}]\cap [\calO_{N}])=m_{F_\calE}.\eeq
\source{very-stable-higgs-bundles-equivariant-multiplicity-mirror-symmetry}
	
	
	
	
	
	
	On the other hand, by the projection formula in $K$-theory \cite[(1)]{piepmeyer-walker} \beq\label{projection}h_*([\calO_{N}] \cap [\calO_{W^+_\calE}])=[\calO_0] \cap h_*([\calO_{W^+_\calE}]) \in K^{\{0\}}_\circ(\calA).\eeq
\source{very-stable-higgs-bundles-equivariant-multiplicity-mirror-symmetry}
	By Lemma~\ref{locfree} $h_*([\calO_{W^+_\calE}])$ is locally free and we  compute $$\chi([\calO_{N}]\cap [\calO_{W^+_\calE}])=\chi(h_*([\calO_{N}]\cap [\calO_{W^+_\calE}]))=\chi([\calO_0]\cap h_*([\calO_{W^+_\calE}]))=\chi([h_*([\calO_{W^+_\calE}])_0])=\rank(h_*(\calO_{W^+_\calE})).$$
	Thus the multiplicity $m_{F_\calE}$ in \eqref{intermult} agrees with $\rank(h_*(\calO_{W^+_\calE}))$ the rank of the locally free sheaf $h_*(\calO_{W^+_\calE})$. 
	
	To determine this rank, we notice that  $W^+_\calE$ is $\T$-invariant. %We will work in $\T$-equivariant $K$-theory \cite[\S 5]{chriss-ginzburg}. If we denote by $[\calO_{W^+_\calE}]^\T\in K^\T_{\circ}(\M)$ the $\T$-equivariant $K$-theory class of $ \calO_{W^+_\calE}$ 
	We  compute the equivariant Euler characteristic of $\calO_{W^+_\calE}$  as \beq \label{one}\chi_\T(\calO_{W^+_\calE})=\chi_\T(\calO_{T_\calE^+})= \chi_\T(\Sym(T_\calE^{+*}))\in\Z((t))\eeq since the $\T$-action on $W^+_\calE$ is modelled on the $\T$-module $T_\calE^+$ i.e. $W^+_\calE\cong T_\calE^+$ are $\T$-equivariantly isomorphic (c.f. \eqref{model}). 
\source{very-stable-higgs-bundles-equivariant-multiplicity-mirror-symmetry}
	
	On the other hand because $h:W^+_{\calE}\to \calA$ is finite the $\T$-equivariant sheaf $h_*(\calO_{W^+_\calE})$ is locally free. By the equivariant  Serre conjecture proved for $\T$ in \cite{equserre} the $\T$-equivariant locally free sheaf $h_*(\calO_{W^+_\calE})$ on the affine space $\calA$ is trivial. In other words \beq\label{hitchinpush}h_*(\calO_{W^+_\calE})\cong h_*(\calO_{W^+_\calE})_0\times \calA,\eeq where the $\T$-action on the $0$ fibre  $h_*(\calO_{W^+_\calE})_0$ is inherited from the $\T$-equivariant structure on $h_*(\calO_{W^+_\calE})$ and the $\T$ action on $\calA$ is the standard one.  Thus we  compute \beq \label{other}\chi_\T(\calO_{W^+_\calE})=\chi_\T(Rh_*(\calO_{W^+_\calE}))=\chi_\T(h_*(\calO_{W^+_\calE}))=\chi_\T(h_*(\calO_{W^+_\calE})_0) \chi_\T(\Sym(\calA^*)) \in \Z((t)).\eeq Here  $\chi_\T(h_*(\calO_{W^+_\calE})_0)$ of \eqref{trunchar} is thought of as the $\T$-equivariant Euler characteristic over the point $0\in \calA$.
\source{very-stable-higgs-bundles-equivariant-multiplicity-mirror-symmetry}
	Comparing \eqref{one} and \eqref{other} we see that the equivariant Euler characteristic of the $\T$-module $h_*(\calO_{W^+_\calE})_0$ can be expressed as \beq \label{t-character}\chi_\T(h_*(\calO_{W^+_\calE})_0)={ \chi_\T(\Sym(T^{+*}_\calE))\over \chi_\T(\Sym(\calA^*))}\in \Z((t)).\eeq In particular, the dimension of $h_*(\calO_{W^+_\calE})_0$ is $$\rank(h_*(\calO_{W^+_\calE}))=\dim(h_*(\calO_{W^+_\calE})_0)=\left. { \chi_\T(\Sym(T^{+*}_\calE))\over \chi_\T(\Sym(\calA^*))}\right|_{t=1}.$$
\source{very-stable-higgs-bundles-equivariant-multiplicity-mirror-symmetry}
	This proves the first statement. For the second statement note that 
	when $\calE\in \M^{s\T}$ is very stable then $W^+_\calE\subset \M$ is closed and for   $a\in\calA$ \beq \label{generic}h_*([\calO_{W^+_\calE}]\cap [h^{-1}(a)])=[h_*(\calO_{W^+_\calE})\cap \calO_a]=[h_*(\calO_{W^+_\calE})_a]=m_{F_\calE}\in K(\{a\})\cong \Z.\eeq When $a\in \calA$ generic then by Lemma~\ref{locfree} the intersection of $W^+_\calE$ and $h^{-1}(a)$ is transversal and thus has cardinality $m_{F_\calE}$ from \eqref{generic}.
\source{very-stable-higgs-bundles-equivariant-multiplicity-mirror-symmetry}
	
	
\end{proof}

\begin{definition}
\source{very-stable-higgs-bundles-equivariant-multiplicity-mirror-symmetry}
\notready\label{equmult} For $\calE\in \M^{s
\source{very-stable-higgs-bundles-equivariant-multiplicity-mirror-symmetry}
		\T}$ define the rational function $$m_{\calE}(t):={{ \chi_\T(\Sym(T^{+*}_\calE))\over \chi_\T(\Sym(\calA^*))}}\in \Z((t)).$$ % To simplify notation sometimes we also write $m_\calE(t)$ for $m_{\calE}(t)$. 
	We call it the {\em virtual equivariant multiplicity} or {\em virtual multiplicity} for short. 
\end{definition}

\begin{corollary}
\source{very-stable-higgs-bundles-equivariant-multiplicity-mirror-symmetry}
\notready\label{eqmult}
\source{very-stable-higgs-bundles-equivariant-multiplicity-mirror-symmetry}
	When $\calE\in \M^{s
		\T}$ is very stable $m_{\calE}(t)$ is a palindromic and monic polynomial with non-negative integer coefficients, such that $m_{\calE}(1)=m_{F_\calE}$ is the multiplicity of the component $N_{F_\calE}\subset N$ in the nilpotent cone.
\end{corollary}
\begin{proof}
	When $\calE$ is very stable  by \eqref{t-character} $m_{\calE}(t)$ is the Laurent polynomial $\chi_\T(h_*(\calO_{W^+_\calE})_0)\in \Z[t,t^{-1}]$, the character of the finite dimensional $\T$-module $h_*(\calO_{W^+_\calE})_0$. In particular, $m_{\calE}(t)$ has non-negative coefficients.  
	
	When all the weights of $\T$ on a finite dimensional $\T$-module $V$ are positive, then ${1/ \chi_{\T}(\Sym(V^*))}$ is a monic Laurent series i.e. $$\left({1\over \chi_{\T}(\Sym(V^*))}\right)_{t=0}=1$$ from \eqref{sym}. It follows that $m_{\calE}(t)\in \Z[t]$ is a polynomial and $m_{\calE}(0)=1$. 
	
	As  $$\prod_{\lambda>0}{( 1- t^{\lambda})^{\dim(V_\lambda)}}=(-1)^{\dim V}t^{\sum \lambda\dim(V_\lambda)} \prod_{\lambda>0}{( 1- t^{-\lambda})^{\dim(V_\lambda)}}$$ we get  \beq\label{palindromic}m_{\calE}(t)=m_{\calE}(t^{-1}) t^{\sum \lambda\left(\dim(\calA_\lambda)-\dim((T^+_\calE)_\lambda)\right)}=m_{\calE}(t^{-1}) t^{\deg(m_{\calE}(t))}
\source{very-stable-higgs-bundles-equivariant-multiplicity-mirror-symmetry}
	\eeq where $$\deg(m_{\calE}(t))=\sum \lambda\left(\dim(\calA_\lambda)-\dim((T^+_\calE)_\lambda)\right)$$ is the degree of the polynomial $m_{\calE}(t)$. Thus $m_{\calE}(t)$ is palindromic.  As $m_{\calE}(0)=1$  it is monic too. The result follows.   
\end{proof}

\begin{remark}
\source{very-stable-higgs-bundles-equivariant-multiplicity-mirror-symmetry}
\notready In Section~\ref{further} we shall see that, in the cases considered there, $m_{\calE}(t^2)$ is a product of Poincar\'e polynomials of compact homogeneous spaces and hence monic and palindromic by Poincar\'e duality. 
\end{remark} 

\begin{remark}
\source{very-stable-higgs-bundles-equivariant-multiplicity-mirror-symmetry}
\notready For a very stable $\calE\in \M^{s\T}$ we have $m_{\calE}(t)=\chi_\T(h_*(\calO_{W^+_{\calE}})_0)$ and so $\T$ acts on the vector space $h_*(\calO_{W^+_{\calE}})_0$ with non-positive weights and the multiplicity of the trivial character of $\T$ is $1$ in $h_*(\calO_{W^+_{\calE}})_0$.
	
\end{remark}

\begin{definition}
\source{very-stable-higgs-bundles-equivariant-multiplicity-mirror-symmetry}
\notready When $\calE\in M^{s\T}$ is very stable we call $m_{\calE}(t)$ the {\em equivariant multiplicity} of $N_{F_\calE}\subset N$. 
\end{definition}
\begin{remark}
\source{very-stable-higgs-bundles-equivariant-multiplicity-mirror-symmetry}
\notready Note that by the $\T$-equivariant analogue of the projection formula \eqref{projection} we see that the equivariant multiplicity at a very stable Higgs bundle satisfies $$m_{\calE}(t)=[\calO_N]^\T\cap [\calO_{W^+_\calE}]^\T\in K_\circ^\T(\{\calE\})\cong \Z[t,t^{-1}]. $$ Hence the name. 
\end{remark}
\begin{remark}
\source{very-stable-higgs-bundles-equivariant-multiplicity-mirror-symmetry}
\notready \label{multN}We know from \cite[Proposition 3.5]{laumon} that   very stable rank $n$ vector bundles $E\in \calN_n$ exist.  Then $\calE=(E,0)$ is a very stable Higgs bundle, and $T^+_\calE\cong T^*_E\calN$ and $\T$ acts with weight one. Thus  \beq \label{cotchar}\chi_\T(\Sym(T^{+*}_\calE))={1\over (1-t)^{n^2(g-1)+1}}\eeq  while $$\chi_\T(\Sym(\calA^*))={1\over (1-t)^g(1-t^2)^{3g-3}(1-t^3)^{5g-5}\cdots (1-t^n)^{(2n-1)(g-1)}}.$$ Thus the equivariant multiplicity in this case
\source{very-stable-higgs-bundles-equivariant-multiplicity-mirror-symmetry}
	is \beq \label{eqmultN}m_{\calE}(t)={{ \chi_\T(\Sym(T^{+*}_\calE))\over \chi_\T(\Sym(\calA^*))}}=[2]_t^{3g-3}[3]_t^{5g-5}\cdots [{n}]_t^{(2n-1)(g-1)}, \eeq  where $$[n]_t:=\frac{1-t^n}{1-t}=1+t+\dots+t^{n-1}$$ are the quantum integers. So by Theorem~\ref{multiplicity} the multiplicity of the component of the nilpotent cone $N$ whose reduced subvariety is isomorphic with $\calN$, the moduli space of rank $n$ stable bundles, is    $$m_{\calN}=2^{3g-3}3^{5g-5}\cdots n^{(2n-1)(g-1)}.$$ This number was also computed in \cite{beauville-etal} as the degree of the dominant map $h^{-1}(a)\to \calN$.   We will see in Subsection~\ref{cotfib} below that this is not a coincidence.  
\source{very-stable-higgs-bundles-equivariant-multiplicity-mirror-symmetry}
	
	Finally, we note that for any other $\calF\in \M^{s\T}$ $${\chi_\T(\Sym(T^{+*}_\calE)) \over \chi_\T(\Sym(T^{+*}_\calF)) }\in \Z[t]$$ is a polynomial from the definition \eqref{sym}, thus we can deduce the following. 
\end{remark}

\begin{corollary}
\source{very-stable-higgs-bundles-equivariant-multiplicity-mirror-symmetry}
\notready  For $\calF\in \M^{s\T}$  very stable, the equivariant multiplicity divides \eqref{eqmultN}:
	$$m_\calF(t)|[2]_t^{3g-3}[3]_t^{5g-5}\cdots [{n}]_t^{(2n-1)(g-1)}.$$  In particular, it is a product of cyclotomic polynomials. 
\end{corollary}

\begin{remark}
\source{very-stable-higgs-bundles-equivariant-multiplicity-mirror-symmetry}
\notready\label{rank2} Let $n=2$. By Corollary~\ref{type111} and Theorem~\ref{laumon} we know that there exists a very stable Higgs bundle $\calE$ in every component $\calN,F_i\in \pi_0(\M_2)$ of the fixed point set of the $\T$-action on $\M$ (see Remark~\ref{remarkn2}).  We see that $T^+_\calE=V_2\oplus V_1$ has only weights $1$ and $2$ and $V_1\cong T^*F_i$, thus $$\chi_\T(\Sym(T^{+*}_\calE))={1\over (1-t)^{\dim F_i}(1-t^2)^{\dim \M/2-\dim F_i}}={1\over(1-t)^{i+g}(1-t^2)^{3g-i-3}}.$$ On the other hand $$\chi_\T(\Sym(\calA^*))={1\over(1-t)^g(1-t^2)^{3g-3}}.$$ Consequently $$m_{\calE}(t)={{ \chi_\T(\Sym(T^{+*}_\calE))\over \chi_\T(\Sym(\calA^*))}}=(1+t)^{i}=[2]_t^i.$$  Thus from Theorem~\ref{multiplicity} we get that $$m_{F_i}=2^i.$$ This result was proved in \cite[Proposition 6]{hitchin1} for $\SL_2$ Higgs bundles ($i$ even) and for twisted $\SL_2$ Higgs bundles ($i$ odd) in \cite{hausel-thaddeus}.	\end{remark}
\source{very-stable-higgs-bundles-equivariant-multiplicity-mirror-symmetry}

\begin{remark}
\source{very-stable-higgs-bundles-equivariant-multiplicity-mirror-symmetry}
\notready \label{counter}
\source{very-stable-higgs-bundles-equivariant-multiplicity-mirror-symmetry}
	One can compute  the rational function ${{ \chi_\T(\Sym(T^{+*}_\calE))/ \chi_\T(\Sym(\calA^*))}}$ for every $\calE\in \M^{s\T}$, which by 
	Corollary~\ref{eqmult} equals the polynomial $m_{\calE}(t)$, when $\calE$ is very stable. This quantity depends only on the ambient $F_\calE\in \pi_0(\M^\T)$ and when it is not a polynomial, we  deduce that there exists no very stable Higgs bundle in $F_\calE$. 
	
	As an example recall \cite{gothen} that a type $(1,2)$ fixed point $\calE$ of the $\T$-action has underlying rank $3$ vector bundle of the form $E=L\oplus V$, where $L$ is a rank $1$ and $V$ is a rank $2$ vector bundle on $C$ and
	the Higgs field only non-trivial in $H^0(C;L^{*}VK)$. If $\ell=\deg(L)$ and $v=\deg(V)$ then $d=\deg(E)=\ell+v$. It is shown in \cite[Proposition 2.5]{gothen} that there is a stable Higgs bundle of this form if $ d/3<\ell<d/3+g-1$ or equivalently $0<2\ell-v<3g-3$. 
	
	We can easily compute the weight spaces of the $\T$-module 
	$T^+_\calE$. It has two weights $1$ and $2$, and we can compute the $2$-weight space isomorphic to $H^0(C;V^*LK)$ whose dimension is $2\ell-v+2g-2$ from Hirzebruch-Riemann-Roch and the fact that $H^1(C;V^*LK)\cong H^0(C;VL^*)^*\cong 0$ from \cite[proof of Proposition 4.2]{gothen}. Because of the homogeneity $1$ symplectic form,  the weight $2$ space is isomorphic with the dual of the weight $-1$ space, thus $2\ell-v+2g-2$ also agrees with half of the Morse index of the ambient fixed point component, matching \cite[Proposition 4.2.i]{gothen}. It follows that  $$\chi_\T(\Sym(T_\calE^{+*}))={1\over(1-t)^{9g-8-(2g-2+2\ell-v)}(1-t^2)^{2\ell-v+2g-2}}.$$ This implies \beq \label{nonpol}m_\calE(t)={{ \chi_\T(\Sym(T^{+*}_\calE))\over \chi_\T(\Sym(\calA^*))}}=(1+t)^{g-1-2\ell+v}(1+t+t^2)^{5g-5}.\eeq
\source{very-stable-higgs-bundles-equivariant-multiplicity-mirror-symmetry}
	This is a polynomial in $t$  if and only if  $2\ell-v\leq g-1$. This shows that the type $(1,2)$ components of $
	\M_3^\T$ where $g-1<2\ell-v<3g-3$ do not contain very stable Higgs bundles.
	
	In fact, we can see this fact directly as follows. Our fixed point is $L\oplus V$ with Higgs field $\phi:L\rightarrow VK$. Assume for simplicity that the rank $3$ bundle has degree $0$. The upward flow consists of extensions $L\rightarrow E\rightarrow V$ where $\Phi$ acting on $L$ and projecting to $V$ is $\phi$.
	
	Let $U\subset V$ be the line bundle generated by $\phi$, or the image of $\phi:LK^*\rightarrow V$. Then 
	$$\deg U\ge \deg L-\deg K=\ell-(2g-2).$$
	Its annihilator is $(V/U)^*\subset V^*$ which has degree 
	$\deg U-\deg V=\deg U+\ell$ so the degree of $L(V/U)^*K$ is 
	$$\deg U+2\ell +(2g-2)\ge 3\ell.$$
	
	If $\ell>(g-1)/3$ then $L(V/U)^*K$ always has a non-trivial section $\psi$ which we can think of as lying in $\Hom(V, LK)$.
	Then $(\phi,\psi):L\oplus V\rightarrow (V\oplus L)\otimes K$ is a Higgs field on $L\oplus V$ and $\langle \phi,\psi\rangle = 0$ means it is nilpotent. 
	
	
\end{remark}

\begin{remark}
\source{very-stable-higgs-bundles-equivariant-multiplicity-mirror-symmetry}
\notready \label{obstruction}  For rank $2$ we know that every component of $\M^\T$ contains a very stable Higgs bundle, thus every component of the nilpotent cone is very stable (see Remark~\ref{rank2}). When a bundle is not very stable, the current term for it is {\em wobbly}. So for  rank $3$ above we found type $(1,2)$ wobbly components of the nilpotent cone by computing $m_{F}(t)$  in \eqref{nonpol} and finding that it is not a polynomial. Interestingly, precisely in the case when $m_{F}(t)$ was not a polynomial we found wobbly directions in $W^+_\calE$.    In fact  for rank $3$ we conjecture that all other type $(1,2)$ components of the nilpotent cone are very stable. As the type $(2,1)$ fixed points behave similarly to the type $(1,2)$ ones by duality, for rank $3$ we can formulate the conjecture that a component  $F\in \pi_0(\M^{s\T})$ contains a very stable Higgs bundle if and only if $m_{\calE}(t)$ (which is independent of $\calE\in F$) is a polynomial for a $\calE\in F$. 
\source{very-stable-higgs-bundles-equivariant-multiplicity-mirror-symmetry}
	
	Adapting the argument for the type $(1,2)$ wobbly fixed point components above to the rank  $4$ case  we find that all stable type $(1,3)$ $\T$-fixed Higgs bundle are  wobbly. However,  the corresponding $m_\calE(t)$ is a polynomial, using the description in \cite[Example 6.6]{garcia-prada-heinloth-schmitt}. 
	Thus the obstruction of integrality of $m_\calE(t)$  for very stable Higgs bundles discussed in Remark  \ref{counter} is not sufficient. In particular, already for rank $4$ we do not have a conjectured list of very stable components of the nilpotent cone. 
	
\end{remark}

\begin{remark}
\source{very-stable-higgs-bundles-equivariant-multiplicity-mirror-symmetry}
\notready\label{remark111}
\source{very-stable-higgs-bundles-equivariant-multiplicity-mirror-symmetry}
	We know from Theorem~\ref{mainverystable} that all type $(1,\dots,1)$ components of the nilpotent cone are very stable and we can explicitly compute  ${{ \chi_\T(\Sym(T^{+*}_\calE))/ \chi_\T(\Sym(\calA^*))}}$ for a fixed point  $\calE\in \M^{s\T}$ of type $(1,\dots,1)$. In this case $E=L_0\oplus \dots \oplus L_{n-1}$ and  $\Phi(L_{i-1})\subset L_{i}K$ is determined by $b_i:=\Phi|_{L_{i-1}}:L_{i-1}\to L_{i}K$ non-zero. 
	We denote $\ell_i=\deg(L_i)$ and $m_i=2g-2-\ell_i+\ell_{i+1}$. As $b_i$ is non-zero $m_i\geq 0$. 
	
	By \eqref{tangent} the tangent space $T_\calE$  to ${\mathcal M}^s$ at this point is the hypercohomology group
	\begin{multline*}T_\calE\cong \H^1\left(\End(E)\stackrel{\ad(\Phi)}{\rightarrow}\End(E)\otimes K\right)\cong \\ \cong \bigoplus_{k=-n+1}^{k=n} \H^1\left(\bigoplus_{i-j=k}\Hom(L_i,L_{j})\stackrel{\ad(\Phi)}{\rightarrow}\bigoplus_{i-j=k-1}\Hom(L_i,L_{j})\otimes K\right)\end{multline*} indexed according to the weights of the $\T$ action on $T_\calE$. Thus we see that \bes T^+_\calE= \bigoplus_{k=1}^{k=n} (T_\calE)_k = \bigoplus_{k=1}^{k=n} \H^1\left(\bigoplus_{i-j=k}\Hom(L_i,L_{j})\stackrel{ad(\Phi)}{\rightarrow}\bigoplus_{i-j=k-1}\Hom(L_i,L_{j})\otimes K\right).\ees Because of the stability of the Higgs bundle $\calE$ \cite[Theorem 4.3]{hausel-vanishing} we get that the zeroth and second hypercohomology of all the complexes vanish, thus \bes\dim T^k_\calE&=&-\sum_{i-j=k} \chi(\Hom(L_i,L_j)) +\sum_{i-j=k-1} \chi(\Hom(L_i,L_j)\otimes K)\\ &=& -\sum_{i-j=k} (-\ell_i+\ell_j+1-g) + \sum_{i-j=k-1} (-\ell_i+\ell_j + g-1) \\ &=& (2n-2k+1)(g-1)-\ell_k+\ell_{n-k+1} \\ &=& (2k-1)(g-1) + m_{k,n-k+1},\ees where $$m_{a,b}=\left\{ \begin{array}{cc}\sum_{a\leq j< b}^b m_j & {\mbox{ when } a<b}\\  0 & {\mbox{ when } a=b} \\  -\sum_{b\leq j<a}^b m_j & {\mbox{ when } a> b} \end{array} \right.$$
	Then a straightforward computation gives \beq\label{explicit}m_\calE(t)={{ \chi_\T(\Sym(T^{+*}_\calE))\over \chi_\T(\Sym(\calA^*))}}={\prod_{j=\lceil{n \over 2}\rceil}^{{n-1}} (1-t^{j+1})^{m_{n-j,j+1}}\over \prod_{j=1}^{\lfloor {n\over 2} \rfloor} (1-t^j)^{m_{j,n-j+1} }}=\prod_{i=1}^{n-1} \left[\begin{array}{c} n \\ i\end{array}\right]^{m_i}_t.\eeq 
\source{very-stable-higgs-bundles-equivariant-multiplicity-mirror-symmetry}
	Here  $$\left[\begin{array}{c} n \\ k \end{array}\right]_t=\prod_{j=1}^k{1-t^{n-j+1} \over 1-t^j }$$ is the quantum binomial coefficient, in particular a polynomial. Incidently, this is the Poincar\'e polynomial of the Grassmanian $\Gr_k(\C^n)$. This coincidence will be further discussed in Subsection~\ref{simple}. 
	
	We thus see that the rational function in \eqref{explicit} is always a polynomial. The polynomiality also follows from the existence of very stable Higgs bundles in each type $(1,1,\dots,1)$  component of $\M^\T$, as proved in Corollary~\ref{type111}. 
\end{remark}

We obtain immediately  
\begin{corollary}
\source{very-stable-higgs-bundles-equivariant-multiplicity-mirror-symmetry}
\notready \label{mult111} For a fixed point component $F_\calE\in \pi_0(\M^T)$  of type $(1,1,...,1)$ as in Remark~\ref{remark111} the multiplicity of $N_{F_\calE}$ in the nilpotent cone $N$ is $$m_{F_\calE}=m_\calE(1)=\prod_{i=1}^{n-1} %\left(\begin{array}{c} n \\ i\end{array}\right)
\source{very-stable-higgs-bundles-equivariant-multiplicity-mirror-symmetry}
	{n \choose i}^{m_i}.$$
\end{corollary}

Another consequence is the following
\begin{corollary}
\source{very-stable-higgs-bundles-equivariant-multiplicity-mirror-symmetry}
\notready For $\calE\in \M^{s\T}_n$ very stable the following are equivalent \begin{enumerate}\item $m_{{F_\calE}}=1$ 		\item $m_{{F_\calE}}(t)=1$  \item $\calE$ is a uniformising Higgs bundle of Remark~\ref{hitchinsection} % i.e. of the form $E=L\oplus LK^{-1}\oplus \dots \oplus LK^{-n+1}$ for some line bundle $L$ with nilpotent Higgs field given by $\Phi|_{LK^{-i+1}}=1: LK^{-i+1}\to LK^{-i}K$. 
		
		\item $W^+_\calE$ is a section of the Hitchin map $h$ 
	\end{enumerate}
\end{corollary}
\begin{proof} Corollary~\ref{eqmult} implies that $m_{\calE}(t)\in \Z_{\geq0}[t]$ is palindromic with non-negative coefficients  showing $1.\Rightarrow 2.$.  By Definition~\ref{equmult}  $m_{\calE}(t)=1$ implies that there are $n$ distinct weights in the $\T$-module $T^+_\calE$, to match those in $\calA$.  This implies in turn that $\calE$ is of type $(1,\dots,1)$. For a type $(1,\dots,1)$ $\calE$ we will have $m_{\calE}(t)=1$ if all $m_i=0$ from \eqref{explicit}. This only happens for a  uniformising Higgs bundle. This shows $2.\Rightarrow 3.$.  From Remark~\ref{hitchinsection} for a uniformising Higgs bundle   $W^+_\calE$ is a section of the Hitchin map. This shows $3.\Rightarrow 4.$. Finally the second part of Theorem~\ref{multiplicity} implies $4.\Rightarrow1.$.
\end{proof}

\subsection{Intersection of upward flows with generic fibres of $h$}
\label{intersection}
\source{very-stable-higgs-bundles-equivariant-multiplicity-mirror-symmetry}



Let ${{a}}=(a_1,\dots,a_n)\in H^0(C;K)\oplus \dots \oplus H^0(C;K^n)=\calA$ be coefficients that define a smooth spectral curve $C_a\subset T^*C$.   Assume that $\div(b)$ for $\calE$ of type $(1,\dots,1)$ avoids the ramification divisor of $\pi_a:C_a\to C$. In other words over any zero $c\in C$ of $b$ we have $|\pi_a^{-1}(c)|$ has $n$ distinct preimages of $c$ in $C_a$.   

We recall that the spectral curve is a subscheme of $T^*C$ defined as $$C_a=\det(\pi^*\Phi-x)^{-1}(0),$$ where $\pi:T^*C\to C$ is the projection and $x\in H^0(T^*C;\pi^*(K))$ is the tautological section.  For $(E,\Phi)\in h^{-1}(a)$ denote by \beq \label{definitionU}U:={\rm coker}(\pi_a^*(E\otimes K^{-1})\stackrel{\pi_a^*\Phi-x}{\to} \pi_a^*(E)).\eeq 
\source{very-stable-higgs-bundles-equivariant-multiplicity-mirror-symmetry}

Then by \cite[\S 5.1]{hitchin-stable} and \cite[\S 3.7]{beauville-etal}  we have the following proposition, which we prove for completeness as the proof is not readily available in the literature.  

\begin{proposition}
\source{very-stable-higgs-bundles-equivariant-multiplicity-mirror-symmetry}
\notready  
	
	\noindent 1. The direct image of the trivial bundle   \beq \label{canonicalBNR} (\pi_a)_*(\calO_{C_a}\stackrel{x}{\to}\pi_a^*(K))\cong \calE_{\calO_C, a}=(E_{\calO_C},\Phi_a)\in W^+_{\calE_0} \eeq   defines  the canonical section  $W^+_{\calE_0}$ of the Hitchin map (see \eqref{uniform}).
\source{very-stable-higgs-bundles-equivariant-multiplicity-mirror-symmetry}
	
	\noindent 2. If a Higgs bundle $(E,\Phi)$ is the direct image of a line bundle $U$:
	\beq \label{utoephi}(\pi_a)_*(U\stackrel{x}{\to} U\otimes \pi_a^*(K))=(E,\Phi).\eeq then  on $C_a$ we have the exact sequence \beq\label{exactu}0\to  U\otimes \pi_a^*(K^{-n}) \to \pi_a^*(E\otimes K^{-1})\stackrel{\pi_a^*\Phi-x}{\to} \pi_a^*(E) \to U\to 0. \eeq 
\source{very-stable-higgs-bundles-equivariant-multiplicity-mirror-symmetry}
\end{proposition}

\begin{proof}  To prove \eqref{canonicalBNR} we note that the pushforward $(\pi_a)_*(\calO_{C_a})$ is the $\calO_C$-algebra ${\rm Sym}(K^{-1})/{\mathcal I}$ where the ideal sheaf ${\mathcal I}$ is generated by the image of $u:K^{-n}\to {\rm Sym}(K^{-1})$ the sum  of $a_i:K^{-n}\to K^{-n+i}$. Thus we see that as a sheaf $(\pi_a)_*(\calO_C)\cong \calE_{\calO_C}$ and the algebra structure on it induces the Higgs field $\Phi_a$ \eqref{companion}  the companion matrix of $a$. 
	
	
	For the proof of the second statement we also denote 
	by $$U^\prime:={\rm ker}(\pi_a^*(E\otimes K^{-1})\stackrel{\pi_a^*\Phi-x}{\to} \pi_a^*(E))$$ the kernel. Then the pushforward of the middle square of
	the diagram $$\begin{array}{ccccccccc} 0\to  &U^\prime& \to &\pi_a^*(E\otimes K^{-1})&\stackrel{\pi_a^*\Phi-x}{\to} &\pi_a^*(E) &\to &U& \to 0 \\  & x \downarrow    &&   x \downarrow  &&  x \downarrow     &&   x \downarrow   & \\ 0\to  &U^\prime\otimes \pi_a^*(K)& \to &\pi_a^*(E)&\stackrel{\pi_a^*\Phi-x}{\to} &\pi_a^*(E\otimes K) &\to &U\otimes \pi_a^*(K) & \to 0 \end{array}$$
	is the middle square of the following diagram 
	
	$$\begin{array}{ccccccccc} 0\to  &EK^{-n}& \stackrel{\Psi}{\to }&E_{\calO_C}\otimes E K^{-1}&\stackrel{Id_{E_{\calO_C}} \otimes \Phi-\Phi_a\otimes Id_{EK^{-1}} }{\longrightarrow} &E_{\calO_C}\otimes E &\stackrel{\Xi}{\to }&E& \to 0 \\  & _\Phi \downarrow    &&   _{\Phi_a\otimes Id_E} \downarrow  &&   _{\Phi_a\otimes Id_E} \downarrow     &&   _\Phi \downarrow   & \\ 0\to  &EK^{-n+1}&  \stackrel{\Psi}{\to} &E_{\calO_C}\otimes E&\stackrel{Id_{E_{\calO_C}} \otimes \Phi-\Phi_a\otimes Id_{EK^{-1}} }{\longrightarrow} &E_{\calO_C}\otimes EK &\stackrel{\Xi}{\to} &E K & \to 0 \end{array}.$$ The rest of the diagram is given by the maps
	$$\Psi=\left(\begin{array}{c} \Psi_1\\ \Psi_2 \\ \vdots \\ \Psi_n \end{array}\right):EK^{-n}\to EK^{-1}\oplus EK^{-2}\oplus \dots \oplus EK^{-n}$$  with  $$\Psi_i=\Phi^{n-i}+a_1 \Phi^{n-i-1}+\dots + a_{n-i} Id_{EK^{-n}}:EK^{-n}\to EK^{-i}$$ and $\Xi=(\Xi_1,\Xi_2,\dots,\Xi_n): E\oplus EK^{-1}\oplus \dots \oplus EK^{-n+1}\to E$ is the map given by $\Xi_i=\Phi^{i-1}:EK^{i-1}\to E$. It is straightforward to check that the second  diagram commutes due to the Cayley-Hamilton identity $0=\Phi^n+a_1\Phi^{n-1}+\dots + a_{n} Id_E$ .  Therefore the second diagram is the pushforward of the first. This implies  \eqref{utoephi} and that $U^\prime\cong U\otimes \pi_a^*(K^{-n})$.  The result follows. 
\end{proof} 



%  We have the following

%\begin{proposition}\label{chaindivisors} Assume that $b$ is non-vanishing at the discriminental locus of the ramified covering $C_a\to C$. Then we have the following
%\begin{enumerate} 
%\item when $(f_i)_{\tilde{c}}=0$ then $(b_{i}\circ\dots \circ b_{n-1})_c=0$ for $\pi_K(\tilde{c})=c$.
%\item when $(b_i)_c=0$ at $c\in C$   then $(f_i)_{\tilde{c}}=0$ precisely for $n-i$ of the lifts $\tilde{c}\in \pi_K^{-1}(\{c\})\subset C_a$
%\end{enumerate}
%	\end{proposition}

%\begin{proof}To prove {\em 1.} let us assume that $(f_i)_{\tilde{c}}=0$ for some $\tilde{c}\in C_a$ such that $\pi_K(\tilde{c})=c\in C$. We can assume that $i$ is the largest with this property, i.e. $(f_{i+1})_{\tilde{c}}\neq 0$. Let $\tilde{c}=(c,\lambda)\in C_a$ where $\lambda\in K_c$ is an eigenvalue of $\Phi_c$. From the two assumptions    ${\rm im}(\Phi_c-\lambda Id_{E_c})\supset (E_iK)_c$ but ${\rm im}(\Phi_c-\lambda Id_{E_c})\not\supset (E_{i+1}K)_c$. Thus ${\rm im}(\Phi_c|_{E_i})\subset  (E_iK)_c$, consequently $(b_i)_c=0$.




%	Then there exists some $w\in E\otimes K^*$ and $\lambda_n\in K$ such that $$v=\Phi(w)-\lambda_n w$$ where $\lambda_n\in K_c$ is an eigenvalue of $\Phi$ at $c:=\pi_{E_1}(v)\in C$, such that $$\pi_{\pi_a^*(E_1)}(\tilde{v})=(c,\lambda_n)\in T^*C\cong K.$$ We compute \begin{eqnarray} \nonumber \Phi^{n-1}(v)&=&\Phi^{n-1}(\Phi(w)-\lambda_n w)\\&=&\Phi^n(w)-\lambda_n(\Phi^{n-1}(w))\nonumber \\&=&e_1\Phi^{n-1}(w)-e_2\Phi^{n-2}(w)+e_3\Phi^{n-3}(w)-\dots (-1)^n e_n w -\lambda_n(\Phi^{n-1}(w))\label{cayley-hamilton}\\&=& e_1^\prime \Phi^{n-2}(\Phi(w)-\lambda_nw)-e^\prime_2\Phi^{n-2}(w)+e_3\Phi^{n-3}(w)-\dots (-1)^n e_n w\nonumber\\&=&
%	e_1^\prime\Phi^{n-2}(v)-e_2^\prime \Phi^{n-3}(v)+e_3^\prime \Phi^{n-4}(v)-\dots (-1)^{n-1}e_{n-1}^\prime v \label{result}
%	\end{eqnarray}
%where $e_i\in K_c^i$ is the $i$th symmetric polynomial in $\lambda_1,\dots,\lambda_n \in K_c$ the eigenvalues of $\Phi_c$, while $e^\prime_i\in K_c^i$ is the $i$th symmetric polynomial in $\lambda_1,\dots,\lambda_{n-1} \in K_c$. In \eqref{cayley-hamilton} we used the Cayley-Hamilton theorem for $\Phi_c$ and in \eqref{result} we used that $e_{i+1}=e_{i+1}^\prime+\lambda_n e_i^\prime$. 

%By  \eqref{preserves} and induction on $i$ we get that  $\Phi^i(v)$ lives in $E_{i+1}K^i$ so we can deduce from \eqref{result} that $\Phi^{n-1}(v)\in E_{n-1}$ and therefore $b(v)=0$ from \eqref{b} as claimed. 

%For the other direction suppose that $(b_i)_c=0$. This means that $$\Phi_c((E_i)_c)\subset (E_iK)_c.$$  Thus $\Phi_c$ leaves $E_i$ invariant. As, by assumption, the zeroes of $b$ avoid the ramification divisor of $\pi_a$ we see that $\Phi_c$ has $n$ distinct eigenvalues $\lambda_1,\dots, \lambda_n\in K_c$. Without loss of generality we can assume that the eigenvectors $w_1,\dots,w_{i}\in E_i$ in the $\Phi_c$-invariant subspace $E_i$ correspond to the first $i\leq n-1$ of these eigenvalues. Let $\lambda_j$ be an eigenvalue for $i<j\leq n$. Then the linear map $$\Phi_c-\lambda_j Id_{(E_i)_c} : (E_{i})_c\to (E_{i}K)_c $$ satisfies $$\Phi_c(w_i)-\lambda_jw_i=(\lambda_i-\lambda_j)w_i$$ and so it is surjective. In particular, for any $v\in (E_i)_c$ there is a $w\in (E_i)_c$ such that $$\Phi_c(w)-\lambda_j w =v.$$ This implies that for $i<j\leq n$ we have $(f_i)_{\tilde{c}_j}=0$ at $\tilde{c}_j=(c,\lambda_j)\in C_a$ . Finally, for $0<j\leq i$ the linear map $\Phi_c-\lambda_j Id_{E_c} $ sends $w_j$ to $0$ and every other eigenvector to a scalar multiple, thus $w_j\in E_i$ is not in the image of $\Phi_c-\lambda_j Id_{E_c} $. Thus $(f_i)_{\tilde{c}_j}\neq 0$ for $0<j\leq i$.  The result follows.		\end{proof}

%\begin{proposition} \label{strict} Assume that ${\rm div}(b)$ avoids the ramification divisor of $\pi_a:C_a\to C$ and for all $i$ the homomorphism of line bundles $b_i\in H^0(M_i^*M_{i+1}K)$ have no repeated zero. Then the cardinality of the intersection satisfies \beq \label{estimate}|W^+_\calE\cap h^{-1}(a)|\leq \prod_{i=1}^{n-1}% \left(\begin{array}{c} n \\ i\end{array}\right) 
%{n\choose i}	^{m_i}.\eeq Equality implies that  $b$ has no repeated  zero. 
%	\end{proposition}


Now let $\calE$ be, as in Remark~\ref{remark111}, a very stable Higgs bundle of type $(1,1,\dots,1)$. We know from Proposition~\ref{filtrationexists} that $(E,\Phi)\in \W^+_{\calE}$ if and only if there exists a filtration with a full flag of subbundles as in \eqref{fullfiltration} which is compatible with the Higgs field \eqref{compatible}. Recall the maps $b_i$ and $b$ from \eqref{bi} and \eqref{be} respectively. 

Assume that $(E,\Phi)\in W^+_\calE\cap h^{-1}(a).$  Denote by $f_i:\pi_a^*(E_i)\to U$ the composition of the  embedding $\pi_a^*(E_i)\subset \pi_a^*(E)$ and projection $\pi_a^*(E)\to U$ arising from the definition \eqref{definitionU} of $U$. Furthermore, denote by $U_i=\im(f_i)$ the image sheaf.  As a non-trivial  subsheaf of the invertible sheaf $U$ it is itself an invertible sheaf. We have thus constructed a filtration $$U_1\subset U_2\subset \dots\subset U_{n-1}\subset U_{n}=U$$ by invertible subsheaves. Denote  the induced map $\tilde{f}_i: U_i\to U$ and by $\Delta_i:=\div(\tilde{f}_i)$ the divisor of zeroes of $\tilde{f}_i\in H^0(C;U_i^*U)$. Note that the divisor $\Delta_1$ determines $$U=\pi_a^*(E_1)(\Delta_1),$$ and so in turn by \eqref{utoephi} it also determines our Higgs bundle $(E,\Phi)\in W^+_\calE\cap h^{-1}(a)$. We get the chain of divisors on $C$ $$\Delta_1\geq \Delta_2\geq \dots \geq \Delta_{n-1}\geq \Delta_n=0.$$ 

\begin{proposition}
\source{very-stable-higgs-bundles-equivariant-multiplicity-mirror-symmetry}
\notready \label{vsintersect} Let $\calE\in \M^{s\T}$ be a type $(1,...,1)$ very stable Higgs bundle. Let $a\in \calA$ be such that the spectral curve $C_a\subset K$ is smooth and $\div(b)$ avoids the ramification divisor of $\pi_a$. \begin{enumerate} \item  If $(E,\Phi)\in W^+_\calE\cap h^{-1}(a)$ then the effective divisor $\Delta_i-\Delta_{i+1}\geq 0$ is reduced and supported on  $(n-i)$ preimages in $C_a$ of $\pi_a$ over each zero of $b_i$ in $C$. More precisely if $c\in C$ is a zero of $b_i$ then $\Phi_c$ leaves $(E_i)_c$ invariant and  $\tilde{c}\in \pi_a^{-1}({c})$ is in the support of $\Delta_i-\Delta_{i+1}$ if and only if $\tilde{c}\in K_c$ is not an eigenvalue for $\Phi_c|_{(E_i)_c}:(E_i)_c\to(E_i)_c K_c $. 
\source{very-stable-higgs-bundles-equivariant-multiplicity-mirror-symmetry}
		
		\item 	On the other hand for any choice of an $(n-i)$ element subset of $\pi_a^{-1}(c)$ for all zeroes $c$ of all the $b_i$, there is a unique $(E,\Phi)\in W_\calE^+\cap h^{-1}(a)$ such that the corresponding invertible sheaf on the spectral curve $C_a$ has the form $U=\pi_a^*(E_1)(\Delta_1)$ for $\Delta_1$ effective and reduced and supported at our chosen points in $C_a$. 
		
		\item	Finally, any two distinct choices of such $(n-i)$ element subsets of $\pi_a^{-1}(c)$ over all zeroes of $b_i$ for all $i$ gives non-isomorphic Higgs bundles. In particular, $$|W^+_\calE\cap h^{-1}(a)|= \prod_{i=1}^{n-1}{n\choose i}	^{m_i},$$ where $\deg(b_i)=m_i$. \end{enumerate}	\end{proposition}

\begin{proof}
	First we prove that $\Delta_i-\Delta_{i+1}\geq 0$ is supported at precisely $(n-i)$ preimages in $C_a$ over each zero of $b_i$ in $C$ with no multiplicity, i.e. $\Delta_i-\Delta_{i+1}$ is reduced.  
	
	Note that $\Phi(E_i)\subset E_{i+1}K$ thus we have the map $$p_i:=\pi_a^*(\Phi)-x:\pi_a^*(E_iK^*)\to \pi_a^*(E_{i+1}).$$ Let $\overline{\im(p_i)}\leq \pi_a^*(E_{i+1})$ denote the subbundle which is the saturation of the subsheaf $\im(p_i)\subset \pi_a^*(E_{i+1})$. Let $V_{i+1}:=\pi_a^*(E_{i+1})/\overline{\im(p_i)}$ be the quotient bundle and $$g_{i+1}:\pi_a^*(E_{i+1})\to V_{i+1}$$ the quotient map.  Another way to think about $V_{i+1}$ is that it is the quotient of the sheaf ${\rm coker}(p_i)$ by its torsion: $$V_{i+1}= \coker(p_i)/T(\coker(p_i)).$$ As $p_i$ is generically an embedding $V_{i+1}$ is a line bundle. We have a map $$h_i:V_{i}\to V_{i+1}.$$ This can be seen by  noting that the embedding $\pi_a^*(E_i)\subset \pi_a^*(E_{i+1})$ induces a map $\coker(p_i)\to \coker{(p_{i+1})}$ which induces a map on the quotients by torsion. By chasing around definitions we see that \beq \begin{array}{ccc} \pi_a^*(E_{i+1}) &\stackrel{g_{i+1}}{\to} &V_{i+1} \label{diagram} \\  \hookuparrow && \ \ \ \uparrow_{\scriptsize h_{i}} \\\pi_a^*(E_{i}) &\stackrel{g_{i}}{\to}  &V_{i}  \end{array}\eeq is commutative.
\source{very-stable-higgs-bundles-equivariant-multiplicity-mirror-symmetry}
	We define $V_1:=\pi_a^*(E_1)$ and $h_1:V_1\to V_2$ as $g_{2}|_{\pi_a^*(E_1)}$ so that the diagram \eqref{diagram} is still commutative for $i=1$. Finally, we have a map $h_{n}:V_{n}\to U$ induced by the map $$\coker(p_{n-1})\to \coker(\pi_a^*(\Phi)-x)=U.$$  This induces the commutative diagram
	
	\beq \label{tricomm}
\source{very-stable-higgs-bundles-equivariant-multiplicity-mirror-symmetry}
	\begin{tikzcd}
	& V_n \arrow{dr}{h_n} \\
	\pi_a^*(E) \arrow{ur}{g_n} \arrow{rr}{f_n} && U
	\end{tikzcd}
	\eeq
	
	If we denote by $$h_{\geq i}:=h_{n}\circ \dots \circ h_i : V_i\to U,$$ then we see from the commutativity of
	the diagrams \eqref{diagram} and \eqref{tricomm} that $f_i=h_{\geq i}\circ g_i$. Because $g_i$ is surjective we can deduce 
	 ${\mathrm{im}} (h_{\geq i})={\mathrm{im}}(f_i)=U_i.$ In particular, $h_{\geq i}:V_i\to U_i$ is surjective and thus an isomorphism. 
	 Thus $\Delta_i-\Delta_{i+1}=\div(h_i)$. 
	
	The map $g_{i+1}|_{\pi_a^*(E_{i})}:\pi_a^*(E_{i})\to V_{i+1}$ vanishes at $\tilde{c}=(c,\lambda)\in C_a$ provided that $$\pi_a^*(E_{i})_{\tilde{c}}\subset p_i(\pi_a^*(E_iK^*)_{\tilde{c}}).$$ This happens when
	$$(E_{i})_c\subset (\Phi-\lambda Id_{E})((E_iK^*)_c).$$ As $(E_{i})_c$ and $(E_iK^*)_c$ are of the same dimension this 
	is equivalent to \beq \label{equal}(E_{i})_c= (\Phi-\lambda Id_{E})((E_iK^*)_c).\eeq  This implies $\Phi((E_i)_c)\subset (E_iK)_c$, i.e. that $(E_i)_c$ is $\Phi$-invariant, which implies that $(b_i)_c=0$. As by assumption $c$ avoids the ramification divisor of $\pi_a$ we see that $\Phi_c$ has distinct eigenvalues, i.e. it is regular semisimple.  Finally, when $\lambda$ is not an eigenvalue of $\Phi|_{(E_i)_c}$ then $(\Phi-\lambda Id_{E})$ is injective so gives an isomorphism between   $(E_iK^*)_c$ and $(E_{i})_c$ and thus \eqref{equal} holds.
\source{very-stable-higgs-bundles-equivariant-multiplicity-mirror-symmetry}
	
	Thus we see that $g_{i+1}|_{\pi_a^*(E_i)}$ vanishes at $n-i$ preimages of each zero of $b_i$, the ones corresponding to the eigenvectors of $\Phi_c$ not in $(E_i)_c$. From the commutativity of \eqref{diagram} and surjectivity of $g_i$ we see that the same follows for $h_i$. 
	
	We  compute \bes \deg(\Delta_1)=\sum_{i=1}^{n-1}\deg(\Delta_i-\Delta_{i+1}) = \sum_{i=1}^{n-1}\deg(h_i)  \geq \sum_{i=1}^{n-1} m_i (n-i)=\deg(\Delta_1).\ees Indeed, by  assumption $\div(b_i)$ is reduced, $\deg(\div(b_i))=m_i$, so $\deg(h_i)\geq m_i(n-i)$  and we have \bes \sum_{i=1}^{n-1}m_i(n-i)=\sum_{i=1}^{n-1}(-\ell_i+\ell_{i+1}+2g-2)(n-i)&=&-n\ell_1 +\sum_{i=1}^{n-1} \ell_i + n(n-1)(g-1)\\&=&-n\deg(E_1)+\deg(E)+\deg(K^{n\choose 2})\\&=&-\deg\pi_a^*(E_1)+\deg(U)=\deg(\Delta_1).\ees Here  we used $\deg(E)=\deg(U)-\deg(K^{{n\choose 2}})$ for which see \cite[\S 4]{beauville-etal}. 
	It follows that $\Delta_i-\Delta_{i-1}\geq 0$ is reduced. Over a zero of $b_i$ the morphism $h_i:V_i\to V_{i+1}$ is zero at precisely those points which correspond to eigenvectors of $\Phi_c$  not in $(E_i)_c$. This proves the first part of the statement. 
	% Thus $\deg(\Delta_i-\Delta_{i-1})=\deg(h_i)=m_i(n-i)$. 
	
	For the second part of Proposition~\ref{vsintersect}, we will construct a Higgs bundle such that $\Delta_1\subset C_a$ is the union of the prescribed choices of an $(n-i)$ tuple in $\pi_a^{-1}(c)$ over every zero $c$ of every $b_i$. The construction is by induction on $\deg(b_i).$ When $\deg(b_i)=0$ then $U=\pi_a^*(E_1)$ is unique and given by the Higgs bundle $\calE_{E_1,a}$ defined   in Remark~\ref{hitchinsection}. 
	
	Recall the notation $\calE_{\delta}$ for a type $(1,\dots,1)$ Higgs bundle from Example~\ref{ex111}. Suppose  that we have constructed a Higgs bundle in $(E,\Phi)\in W^+_{\calE_{\d}}\cap h^{-1}(a)$ for any choice of an $(n-i)$ tuple in $\pi_a^{-1}(c)$ over every zero $c$ of every $b_i$. Let $c\in C$ be a point which avoids the ramification divisor of $\pi_a$ and $b(c)\neq 0$. Denote by $b^\prime_j:=b_j$ and $\delta_j^\prime:=\div(b_j)$ except when $j=i$,  when we let $b^\prime_i:=b_i s_c$ or equivalently $\d^\prime_i:=\d_i+c$. Assume that $\calE_{\d^\prime}$ is still stable.  Then $\Phi_c$ is regular semisimple and we can take an $i$-dimensional $\Phi_c$-invariant subspace $V_i\subset E_c$ corresponding to any $i$ element subset of $\pi_a^{-1}(c)$. Then by Proposition~\ref{modifiedfilt} the Hecke transform  $\calH_{V_i}(E,\Phi)=(E^\prime,\Phi^\prime)$ will be contained in $W^+_{\calE_{\d^\prime}}\cap h^{-1}(a)$, and $b_i(c)=0$ so that $\Phi^\prime_c$ leaves $(E_i)_c$ invariant. By induction the second part of the proposition follows.   
	
	For the last part we note that there are ${n\choose n-i}^{m_i}={n\choose i}^{m_i}$ choices for $\Delta_i-\Delta_{i-1}$. 	This gives us $\prod_{i=1}^{n-1}{n \choose i}^{m_i}$ possibilities for $\Delta_1$.  In theory, two distinct divisors could represent the same line bundle, but in our case, this cannot happen. Namely, a Higgs bundle $(E,\Phi) \in W^+_{\calE_{\d}}\cap h^{-1}(a)$ will have a unique filtration yielding the limiting Higgs bundle $\calE_{\d}$ from Proposition~\ref{filtrationexists}. Then in turn the  filtration  defines a unique subset $\Delta_1\subset C_a$ as above in the proof of part one of the proposition. The result follows.
	
	
	
	%	However, when $b$ has repeated zero, for example for $i_1>i_2$ $b_{i_1}$ and $b_{i_2}$ have a common zero, then over it $\Delta_{i_1}-\Delta_{i_1+1}$ should be contained in $\Delta_{i_2}-\Delta_{i_2+1}$, corresponding to eigenvectors of $\Phi_c$ not in $(E_{i_1})_c$ and not in $(E_{i_2})_c\subset (E_{i_1})_c$ respectively. This reduces the number of possibilities for $\Delta_1$, giving the strict inequality in \eqref{estimate}.
	
\end{proof}

\begin{remark}
\source{very-stable-higgs-bundles-equivariant-multiplicity-mirror-symmetry}
\notreadyWhen $\calE\in \M^{s\T}$ is very stable then $W^+_\calE\subset \M$ is closed and we have \eqref{generic}. %$$h_*([\calO_{W^+_\calE}]\cap [h^{-1}(a)])=[h_*(\calO_{W^+_\calE})\cap \calO_a]=[h_*(\calO_{W^+_\calE})_a]=m_\calE\in K(\{a\})\cong \Z.$$ 
	When $\calE$ is of type $(1,\dots,1)$ from Corollary~\ref{mult111} and Proposition~\ref{vsintersect} we get that $|W^+_\calE\cap h^{-1}(a)|=m_\calE$, provided that $C_a$ is smooth and $\div(b)$ avoids the ramification divisor of $\pi_a$. Consequently, in this case $W^+_\calE$ and $h^{-1}(a)$ intersect transversally. 
	
	Similarly as above in Proposition~\ref{vsintersect} one can prove that $|W^+_\calE\cap h^{-1}(a)|<m_\calE$ when $b$ has repeated $0$'s. The second part of Theorem~\ref{multiplicity}  implies that such type $(1,\dots,1)$ Higgs bundles must be  wobbly. This gives an alternative proof of the ``only if " part of Theorem~\ref{mainverystable}.
\end{remark}

%\begin{theorem}\label{repeated} When $(E^\prime,\Phi^\prime)\in \M^{s\T}$ is of type $(1,1,\dots,1)$ given by the chain $(b_1,\dots,b_{n-1})$ as in \eqref{bi} then we have the following \begin{enumerate} \item
%		when $b=b_{n-1}\circ\dots\circ b_1$ has no repeated zero then we have an equality in \eqref{estimate} for a generic $a\in \calA$ \item when $b$ has repeated zeroes then $(E,\Phi)$ is not very stable. 
%		\end{enumerate}
%	\end{theorem}

%\begin{proof} When $\calE\in \M^{s\T}$ is very stable then $W^+_\calE\subset \M$ is closed and $$h_*([\calO_{W^+_\calE}]\cap [h^{-1}(a)])=[h_*(\calO_{W^+_\calE})\cap \calO_a]=[h_*(\calO_{W^+_\calE})_a]=m_\calE\in K(\{a\})\cong \Z.$$ When $a$ is generic all intersection points in $h^{-1}(a)\cap W^+_\calE$ are transversal as $h|_{W^+_\calE}:W^+_\calE\to \calA$ is generically \'etale by Lemma~\ref{locfree} and because $h^{-1}(a)\subset \M$ is a regular subscheme.  Thus for a very stable Higgs bundle and generic $a\in\calA$ we have \beq \label{multgen}m_\calE= |h^{-1}(a)\cap W^+_\calE|.\eeq 
%	Thus when $\calE\in \M^{s\T}$ is  of type $(1,1,\dots,1)$ with $\div(b)$ reduced then it is very stable by Corollary~\ref{verystable111} and \eqref{multgen} and \eqref{mult111}   implies {\em 1.}

%	For {\em 2.} assume that $b$ has repeated zeroes. If any of the $b_i$'s have repeated zeroes than we can argue as in the proof of Theorem~\ref{verystablerank2thm}.
%	We can thus further suppose that the $b_i$'s have individually no repeated zeroes but their product $b$ does. Then by Proposition~\ref{estimate} $$|W^+_\calE\cap h^{-1}(a)| < \prod_{i=1}^{n-1}{n\choose i}	^{m_i}$$ but 
%	$$m_\calE(1)=\prod_{i=1}^{n-1}{n\choose i}	^{m_i}$$ by \eqref{explicit}. Thus \eqref{multgen} implies that $(E^\prime,\Phi^\prime)$ cannot be very stable in this case. 

%	The result follows. 
%	\end{proof}
%
